% Chapter 2

\chapter{Related Work} % Main chapter title
\label{related_work}

This chapter positions the study against three lines of work: datasets and systems for lay summarisation of biomedical text, the evaluation of faithfulness in abstractive summarisation, and decoding pathologies of neural text generators. It ends with the models and metrics this study builds on.

\section{Lay Summarisation of Biomedical Text}

Automatic lay summarisation has mostly been studied on scientific articles and reviews. \textcite{guo2021lay} introduced the task of generating lay-language summaries of biomedical scientific reviews, released a dataset for it, and identified three sources of difficulty: condensing the key points, explaining background knowledge, and simplifying technical language. \textcite{devaraj2021paragraph} built a parallel corpus of technical and lay summaries of clinical evidence for paragraph-level simplification, proposed a readability measure based on a masked language model, and trained Transformer models that discourage jargon. \textcite{goldsack2022making} released two corpora, PLOS and eLife, that pair biomedical journal articles with plain-language summaries written for the public, and benchmarked standard summarisation methods on them. These corpora were the basis of the BioLaySumm 2023 shared task, in which 20 teams submitted systems, one track of which also required control of readability \parencite{goldsack2023biolaysumm}.

Our setting differs from this work in three ways. First, the source is a sponsor's clinical trial results report, not a journal article or review. Second, the target is not one free-form summary but a sequence of answers to fixed patient questions (\emph{Why was the research needed?}, \emph{Who was in this clinical trial?}, and so on), each paired with mapped sections of the report. Third, we restrict ourselves to models that can be trained and run on a CPU, so that the results describe what is feasible without specialised hardware.

\section{Faithfulness of Abstractive Summaries}

Abstractive summarisers can produce text that is fluent but not supported by the source. \textcite{maynez2020faithfulness} annotated summaries from several neural systems, found unsupported (hallucinated) content in the outputs of every system examined, and showed that textual-entailment scores agree with human faithfulness judgements better than ROUGE does. Two families of automatic faithfulness metrics followed. Question-answering based metrics such as QAGS generate questions from the summary and compare the answers obtained from the summary and from the source \parencite{wang2020qags}. Entailment based metrics such as SummaC split the document and summary into sentences and aggregate sentence-level natural-language-inference scores \parencite{laban2022summac}.

Plain-language summarisation makes this problem sharper, because the target is meant to add explanation that the source does not contain verbatim. A metric that rewards similarity to a human reference cannot by itself distinguish a helpful explanation from an invented fact. This study does not run QA- or NLI-based metrics; it uses simple string-level proxies (Section~\ref{sec:metrics}) and leaves those metrics to future work.

\section{Decoding and Text Degeneration}

Maximisation-based decoding (greedy or beam search) from neural language models tends to produce repetitive loops \parencite{holtzman2020degeneration}. Common remedies include sampling-based decoding, which is not reproducible run to run, and hard constraints such as forbidding any repeated $n$-gram. \textcite{su2022contrastive} proposed contrastive search, which penalises candidate tokens whose representation is too similar to the preceding context. Chapter~\ref{results} shows that, in our setting, the choice between plain greedy decoding and greedy decoding with a no-repeat-trigram constraint changes both the outputs and the conclusions drawn from the metrics.

\section{Models and Metrics Used in This Study}

\paragraph{Models.} T5 casts every task as text-to-text generation with an encoder-decoder Transformer \parencite{raffel2020t5}. Flan-T5 is T5 further fine-tuned on a large collection of tasks phrased as instructions \parencite{chung2024flan}, which makes it usable zero-shot and a natural starting point for fine-tuning. Low-Rank Adaptation (LoRA) freezes the pretrained weights and trains low-rank update matrices inside the Transformer layers \parencite{hu2022lora}, which reduces the memory needed for fine-tuning. DistilBERT is a smaller BERT model trained by knowledge distillation \parencite{sanh2019distilbert}; we use a checkpoint fine-tuned for extractive question answering on SQuAD \parencite{rajpurkar2016squad}. Qwen2.5 is a family of decoder-only instruction-tuned language models \parencite{qwen2024qwen25}; we use its 1.5-billion-parameter instruct variant as a CPU-feasible stand-in for the Mistral-7B-Instruct model \parencite{jiang2023mistral} used in the original thesis. All models are run through the Hugging Face Transformers library \parencite{wolf2020transformers}.

\paragraph{Metrics.} ROUGE measures $n$-gram and longest-common-subsequence overlap with a reference summary \parencite{lin2004rouge}. BLEU measures modified $n$-gram precision with a brevity penalty \parencite{papineni2002bleu}; at the sentence level it needs smoothing, for which we use method 4 of \textcite{chen2014smoothing}. Token-level F1 follows the SQuAD evaluation \parencite{rajpurkar2016squad}. BERTScore matches tokens of the candidate and the reference by the cosine similarity of their contextual embeddings \parencite{zhang2020bertscore}. Readability is measured with the Flesch Reading Ease score \parencite{flesch1948} and the Flesch--Kincaid grade level \parencite{kincaid1975}. Both formulas depend only on sentence length and word length in syllables; they say nothing about whether a text is correct or whether its terms are explained.
